%% PREAMBLE: Le chapitre utilise \defin, dont la définition appelle
%% \index.  Si ce n'est pas déjà fait dans le document maître, ajouter
%% \usepackage{makeidx} puis \makeindex dans le préambule.
%% PREAMBLE: Pour compiler l'éventuel index, utiliser par exemple
%% texindy -C utf8 -L french notes-inf110.idx.

\section{Généralités sur le typage}

\review{Ce chapitre constitue un premier brouillon produit avec l'aide
d'un modèle d'IA. Il doit être relu et vérifié par l'enseignant,
notamment en ce qui concerne les formulations mathématiques, les
preuves esquissées et les choix éditoriaux.}

\subsection{Rôle et variétés des systèmes de typage}

\thingy Informellement, un \defin{système de typage} est une façon
d'affecter aux expressions et aux valeurs manipulées par un langage
informatique un \defin{type} qui contraint leur usage.  Parmi les types
les plus courants, on trouve par exemple les entiers, les booléens, les
chaînes de caractères et les fonctions.

Cette manière de distinguer différentes sortes de données est, dans
une certaine mesure, opposée à la philosophie du codage de Gödel
rencontrée en calculabilité : dans ce dernier, toute donnée finie est
représentée par un entier, tandis qu'un système de typage cherche au
contraire à conserver et exploiter les distinctions entre les
différentes natures de données.

Les buts d'un système de typage sont multiples :
\begin{itemize}
\item détecter le plus tôt possible certaines erreurs de
  programmation : ajouter un entier et une chaîne de caractères est
  probablement une erreur, ou demande au moins une conversion
  explicite ;
\item éviter des plantages ou des problèmes de sécurité, par exemple
  lorsqu'un programme tente d'exécuter comme du code une donnée qui
  n'est qu'un entier ;
\item garantir certaines propriétés du comportement des programmes,
  dans les limites imposées par les résultats d'indécidabilité tels
  que le théorème de Rice ;
\item fournir des informations utiles à l'optimisation, par exemple en
  indiquant qu'une fonction est pure, c'est-à-dire sans effet de bord.
\end{itemize}
L'analogie avec les systèmes d'unités en physique est utile : de même
que l'homogénéité interdit d'ajouter directement une longueur à une
durée, le typage interdit certaines combinaisons d'expressions qui
n'ont pas de sens dans le langage considéré.

\thingy Il existe presque autant de systèmes de typage que de langages
de programmation. Plusieurs distinctions sont particulièrement
importantes.

Un typage est dit \defin{statique} lorsqu'il est vérifié au moment de la
compilation, et \defin{dynamique} lorsqu'il est vérifié pendant
l'exécution. Certains langages combinent les deux approches ; on parle
notamment de typage graduel. Du point de vue du programmeur, on préfère
en général une erreur de compilation à une erreur à l'exécution, et une
erreur à l'exécution à un plantage ou à un comportement indéfini.

Le typage peut être \defin{annoté}, lorsque le programmeur écrit
explicitement les types attendus, ou \defin{inféré}, lorsque le langage
les détermine à partir de l'expression. Cette distinction peut
s'interpréter de la manière suivante : le type est-il une promesse
faite par le langage au programmeur, ou une promesse faite par le
programmeur au langage ?

Il faut également distinguer les systèmes de typage dont les garanties
sont difficiles à contourner de ceux qui permettent des conversions
explicites (\textit{casts}), la manipulation directe des
représentations en mémoire, ou le report de certaines vérifications à
l'exécution.

Enfin, le typage peut être superficiel, par exemple en affirmant
seulement qu'une valeur est une liste, ou plus précis, en affirmant
qu'elle est une liste d'entiers. Il peut aussi enregistrer des
informations qui ne décrivent pas directement la représentation d'une
valeur : pureté d'une fonction, exceptions qu'elle peut soulever,
mutabilité, droits d'accès, utilisation de la mémoire, etc.

\thingy Dans certains langages, les types eux-mêmes sont des objets
manipulables par les programmes : on dit alors qu'ils sont des
\defin{citoyens de première classe}. Cette possibilité pose aussitôt la
question du type des types eux-mêmes. Elle conduit à des systèmes très
expressifs, mais aussi à des difficultés logiques et algorithmiques
importantes.

\subsection{Constructions usuelles sur les types}

\thingy En plus de types de base tels que $\texttt{Nat}$, pour les
entiers naturels, ou $\texttt{Bool}$, pour les booléens, les systèmes de
typage proposent souvent plusieurs opérations de construction de
types.

Les \defin{types produits} représentent les couples, triplets et, plus
généralement, les $k$-uplets. Par exemple,
\[
\texttt{Nat}\times\texttt{Bool}
\]
est le type des paires dont la première composante est un entier et la
seconde un booléen. Lorsque les composantes sont nommées, on obtient
des structures, aussi appelées enregistrements. Le produit vide est le
type trivial $\texttt{Unit}$, qui ne possède qu'une seule valeur.

Les \defin{types sommes} représentent un choix entre plusieurs cas. Le
type
\[
\texttt{Nat}+\texttt{Bool}
\]
contient soit un entier, soit un booléen, avec un sélecteur indiquant
lequel des deux cas est utilisé. Une somme de plusieurs copies du type
$\texttt{Unit}$ donne un type d'énumération, dont la seule information
est précisément ce sélecteur. La somme vide est un type inhabité,
c'est-à-dire un type ne possédant aucune valeur.

Les \defin{types fonctions}, également appelés types exponentiels, sont
notés à l'aide d'une flèche. Ainsi,
\[
\texttt{Nat}\rightarrow\texttt{Bool}
\]
est le type des fonctions de $\texttt{Nat}$ vers $\texttt{Bool}$.

Enfin, on rencontre fréquemment des types de listes. Par exemple,
$\texttt{List}~\texttt{Nat}$ est le type des listes d'entiers.

\thingy Plusieurs fonctionnalités peuvent enrichir ces constructions.

Le \defin{sous-typage} signifie que toute valeur d'un certain type peut
automatiquement être utilisée comme une valeur d'un autre type.

Le \defin{polymorphisme} permet à une expression d'être utilisée avec
plusieurs types. L'exemple fondamental est la fonction identité, dont
on peut écrire informellement le type
\[
(\forall\mathtt{t})\; \mathtt{t}\rightarrow\mathtt{t}.
\]

Une \defin{famille de types} fabrique un type à partir d'un ou plusieurs
autres types. Ainsi, $\texttt{List}$ transforme le type $\mathtt{t}$ en
le type $\texttt{List}~\mathtt{t}$ des listes d'éléments de type
$\mathtt{t}$.

Un \defin{type récursif} est construit, à l'aide des opérations sur les
types, en faisant intervenir le type défini lui-même. Les diapositives
donnent comme exemple schématique
\[
\texttt{Tree}=\texttt{List}~\texttt{Tree}.
\]
\review{Cet exemple décrit des arbres dont chaque nœud est essentiellement
une liste de sous-arbres, mais ne fournit pas de cas de base explicite ;
il convient de vérifier si cette présentation volontairement
schématique doit être conservée.}

Un \defin{type dépendant} est un type qui dépend d'une valeur. Par
exemple, à un entier $k$ on peut associer le type $\texttt{Nat}^k$ des
$k$-uplets d'entiers.

Enfin, un \defin{type opaque}, abstrait ou privé, est un type dont la
représentation interne est cachée. Ses valeurs ne peuvent alors être
manipulées qu'au moyen d'une interface publique.

\subsection{Polymorphisme et tâches d'un système de typage}

\thingy On distingue principalement deux sortes de polymorphisme.

Le \defin{polymorphisme paramétrique}, ou polymorphisme générique,
signifie que la même expression s'applique de la même manière à des
données de types différents. Quelques exemples sont :
\[
\texttt{head} :
(\forall\mathtt{t})\;
\texttt{List}~\mathtt{t}\to\mathtt{t},
\]
\[
\lambda xy.\langle x,y\rangle :
(\forall\mathtt{u},\mathtt{v})\;
\mathtt{u}\to\mathtt{v}\to
\mathtt{u}\times\mathtt{v},
\]
et
\[
\lambda xyz.xz(yz) :
(\forall\mathtt{u},\mathtt{v},\mathtt{w})\;
(\mathtt{u}\to\mathtt{v}\to\mathtt{w})
\to(\mathtt{u}\to\mathtt{v})
\to\mathtt{u}\to\mathtt{w}.
\]
Le polymorphisme ne concerne pas seulement les fonctions. La liste
vide, par exemple, a informellement le type
\[
\texttt{[]} :
(\forall\mathtt{t})\;\texttt{List}~\mathtt{t}.
\]
Dans un langage autorisant les boucles infinies, une expression qui ne
termine jamais peut également recevoir n'importe quel type, puisqu'elle
ne produit jamais de valeur contredisant ce type.

Le \defin{polymorphisme ad hoc}, ou surcharge, désigne au contraire une
opération qui agit différemment selon le type de son argument, ce type
étant connu au moment où l'opération est choisie.

Les frontières terminologiques sont parfois floues : le sous-typage,
et quelquefois même les coercions, sont présentés comme des formes de
polymorphisme.

\thingy Un système de typage peut accomplir plusieurs tâches.

La \defin{vérification de type} consiste à vérifier qu'une expression
annotée possède bien le type annoncé par ses annotations.

L'\defin{inférence de type}, également appelée reconstruction ou
assignation de type, consiste à déterminer le type d'une expression en
l'absence d'annotations, ou en présence d'annotations seulement
partielles. L'algorithme de Hindley--Milner, étudié à la fin de ce
chapitre, est un algorithme fondamental d'inférence pour les langages
fonctionnels à polymorphisme paramétrique.

Dans les systèmes suffisamment complexes, l'inférence, et parfois même
la vérification, peuvent devenir indécidables. Cela se produit
notamment lorsque les types sont des objets de première classe ou
peuvent dépendre de valeurs arbitraires calculées pendant l'exécution.

Enfin, le problème de l'\defin{habitation d'un type} consiste à trouver
un terme ayant un type donné. Cette question est essentielle en
logique, où elle correspondra à la recherche d'une preuve. Par exemple,
on peut demander s'il existe un terme de type
\[
(\forall\mathtt{p},\mathtt{q})\;
((\mathtt{p}\to\mathtt{q})\to\mathtt{p})\to\mathtt{p}.
\]
Dans les langages usuels qui autorisent la non-terminaison, tout type
peut être artificiellement habité par une boucle infinie, y compris un
type censé être vide. Cette observation explique pourquoi la
terminaison sera indispensable lorsque les types seront interprétés
comme des propositions.

\subsection{Autres usages et exemples}

\thingy Les types peuvent annoter des propriétés qui ne sont pas
simplement la nature des valeurs stockées. Ils peuvent indiquer les
exceptions susceptibles d'être soulevées, comme en Java, ou distinguer
les valeurs immuables des références mutables. Ainsi,
$\texttt{Nat}$ peut désigner un entier immuable, tandis que
$\texttt{Ref~Nat}$ désigne une référence mutable vers un entier.

Ils peuvent également enregistrer les effets de bord. En Haskell,
$\texttt{Char}$ représente une valeur pure de type caractère, tandis
que $\texttt{IO~Char}$ représente une action ayant des effets de bord
et renvoyant un caractère.

Le \defin{typage linéaire}, ou les types à unicité, imposent qu'une
valeur soit utilisée une et une seule fois au cours d'un calcul, sauf
opération spéciale. Ils interdisent donc aussi bien sa duplication que
sa perte. Cette information peut servir à la gestion de la mémoire,
comme en Rust, ou à l'annotation des effets de bord, comme en Clean.

\thingy Donnons quelques exemples de systèmes de typage, sans employer
les expressions ambiguës de typage « faible » ou « fort ».

L'assembleur ne possède essentiellement aucun typage : toute donnée
est représentée par une suite de bits. La même remarque vaut, dans des
cadres théoriques différents, pour les machines de Turing, les
fonctions générales récursives et le $\lambda$-calcul non typé.

Le langage C possède un typage annoté et vérifié à la compilation, mais
celui-ci est contournable, notamment à l'aide des pointeurs génériques
de type \texttt{void*}. Python, JavaScript, Perl ou Scheme effectuent
principalement les vérifications à l'exécution. Java combine des
vérifications à la compilation et à l'exécution ; les types génériques
ont été ajoutés avec Java 5. Rust fait interagir le typage et la gestion
de la mémoire au moyen de mécanismes proches du typage affine ou
linéaire.

Les langages fonctionnels ont souvent des systèmes plus complexes.
OCaml dispose notamment de l'inférence de type à la Hindley--Milner, du
polymorphisme paramétrique, de types récursifs et d'un système de
modules. Haskell possède des fonctionnalités comparables, auxquelles
s'ajoutent les classes de types et l'encapsulation des effets de bord
dans des monades. Mercury combine des idées issues de Haskell et de
Prolog. Idris est à la fois un langage fonctionnel et un assistant de
preuve.

Parmi les exemples plus spécialisés, F\# peut intégrer les unités de
mesure au système de typage, Rust peut exclure certaines situations de
concurrence dangereuses, et divers langages garantissent des propriétés
telles que des bornes de complexité, la validation de documents XML ou
l'absence de fuite d'information.

\subsection{Typage, terminaison et correspondance de Curry--Howard}

\thingy\label{typing-and-termination} Un système de typage peut garantir
que tout programme bien typé termine. Ce n'est pas le cas des langages
usuels tels qu'OCaml ou Haskell, dans lesquels un programme bien typé
peut effectuer une boucle infinie.

Si le typage doit lui-même être décidable, garantir la terminaison
impose nécessairement des limites à l'expressivité du langage, en
raison de l'indécidabilité du problème de l'arrêt. Un tel langage ne
peut notamment pas écrire son propre interpréteur universel. Il ne
peut pas non plus autoriser des boucles illimitées, des appels récursifs
non contraints ou un type $\mathtt{t}$ satisfaisant
\[
\mathtt{t}\cong(\mathtt{t}\rightarrow\mathtt{t}).
\]
Son type vide doit réellement être inhabité.

Ces restrictions n'empêchent cependant pas le langage de définir des
fonctions beaucoup plus générales que les seules fonctions primitives
récursives. Rocq, anciennement appelé Coq, fournit un exemple de
langage et d'assistant de preuve garantissant la terminaison.

\thingy La \defin{correspondance de Curry--Howard} établit une analogie,
et dans des systèmes précisément définis une correspondance exacte,
entre :
\begin{itemize}
\item les types et leurs termes, c'est-à-dire les programmes possédant
  ces types ;
\item les propositions et leurs preuves.
\end{itemize}
Par exemple, la preuve évidente de
\[
(A\land B)\Rightarrow A
\]
correspond à la première projection, de type
\[
a\times b\rightarrow a.
\]
De même, le terme $\lambda xy.x$, de type $u\to v\to u$, correspond à
la preuve de la proposition
\[
U\Rightarrow V\Rightarrow U.
\]

Cette correspondance exige certaines précautions. Le langage de
programmation doit garantir la terminaison, faute de quoi une boucle
infinie jouerait le rôle d'une preuve circulaire. La logique
correspondante est en général intuitionniste, donc dépourvue du tiers
exclu. Enfin, les quantificateurs et les types dépendants font
intervenir des subtilités supplémentaires.

Dans un premier aperçu, la correspondance associe :
\begin{itemize}
\item le type fonction $A\to B$ à l'implication $A\Rightarrow B$ ;
\item l'application d'une fonction au \textit{modus ponens} ;
\item l'abstraction à l'ouverture puis à la décharge d'une hypothèse ;
\item le produit $A\times B$ à la conjonction $A\land B$ ;
\item la somme $A+B$ à la disjonction $A\lor B$ ;
\item les types trivial et vide à $\top$ et $\bot$.
\end{itemize}
La négation demandera une discussion particulière.

\thingy Le paradoxe de Curry illustre le danger des preuves
circulaires. Supposons que l'on puisse construire une proposition $A$
telle que
\[
A\Leftrightarrow(A\Rightarrow B).
\]
Alors $B$ est effectivement démontrable : sous l'hypothèse $A$,
l'équivalence donne $A\Rightarrow B$, donc $B$. Cela établit
$A\Rightarrow B$, qui, par l'équivalence, donne à son tour $A$, puis
finalement $B$.

L'erreur ne se trouve donc pas dans cette déduction, mais dans la
supposition tacite qu'une proposition $A$ satisfaisant une telle
auto-référence peut toujours être construite. Une construction de
Quine permet une auto-référence syntaxique, mais ne suffit pas à
obtenir une auto-référence sur la vérité.

\thingy Les \defin{théories des types} utilisées en logique sont des
langages permettant d'écrire à la fois des programmes et des preuves.
Elles peuvent être étudiées comme des systèmes abstraits ou
implémentées dans des assistants de preuve tels que Rocq, Agda ou Lean.
Elles garantissent la terminaison des programmes, ou plus précisément,
dans le cadre du $\lambda$-calcul, la normalisation forte.

Parmi les grandes familles, on peut citer les théories des types à la
Martin-Löf, les variantes du calcul des constructions, et la théorie
homotopique des types. Le $\lambda$-cube de Barendregt organise huit
théories selon la présence ou l'absence de trois fonctionnalités :
\begin{itemize}
\item les termes peuvent dépendre de types ;
\item les types peuvent dépendre de types ;
\item les types peuvent dépendre de termes.
\end{itemize}
Le $\lambda$-calcul simplement typé ne possède aucune de ces trois
fonctionnalités, tandis que le calcul des constructions les possède
toutes.

\section{Le \texorpdfstring{$\lambda$}{lambda}-calcul simplement typé}

\subsection{Syntaxe et règles de typage}

\thingy Le \defin{$\lambda$-calcul simplement typé}, abrégé
$\lambda$CST et également noté $\lambda_\rightarrow$, est une variante
typée du $\lambda$-calcul. Il ne possède qu'une seule opération sur les
types : si $\sigma$ et $\tau$ sont deux types, alors
$\sigma\to\tau$ est le type des fonctions qui prennent un argument de
type $\sigma$ et renvoient une valeur de type $\tau$.

L'application et l'abstraction doivent respecter les types :
\begin{itemize}
\item si $P$ a pour type $\sigma\to\tau$ et $Q$ a pour type $\sigma$,
  alors $PQ$ a pour type $\tau$ ;
\item si $E$ a pour type $\tau$ sous l'hypothèse que la variable libre
  $v$ a pour type $\sigma$, alors
  $\lambda(v:\sigma).E$ a pour type $\sigma\to\tau$.
\end{itemize}

Le typage considéré est annoté, ou « à la Church » : on écrit
$\lambda(v:\sigma).E$, et non simplement $\lambda v.E$.

\thingy Un \defin{type} du $\lambda$CST est défini inductivement par les
deux clauses suivantes :
\begin{itemize}
\item toute variable de type $\alpha,\beta,\gamma,\ldots$ est un type ;
\item si $\sigma$ et $\tau$ sont des types, alors
  $(\sigma\to\tau)$ est un type.
\end{itemize}

Un \defin{préterme} est défini inductivement par :
\begin{itemize}
\item toute variable de terme est un préterme ;
\item si $P$ et $Q$ sont des prétermes, alors $(PQ)$ est un préterme ;
\item si $v$ est une variable, $\sigma$ un type et $E$ un préterme,
  alors $\lambda(v:\sigma).E$ est un préterme.
\end{itemize}

On adopte les conventions de parenthésage
\[
\rho\to\sigma\to\tau
  :=\rho\to(\sigma\to\tau)
\qquad\text{et}\qquad
xyz:=((xy)z).
\]
L'application est donc associative à gauche dans la notation, tandis
que le constructeur de type fonction est associatif à droite.
L'abstraction est moins prioritaire que l'application. On peut abréger
\[
\lambda(x:\sigma).\lambda(t:\tau).E
\]
en $\lambda(x:\sigma,t:\tau).E$ ou
$\lambda(x:\sigma)(t:\tau).E$. Enfin, les termes sont considérés à
renommage près de leurs variables liées, c'est-à-dire à
$\alpha$-conversion près.

\begin{defn}
Un \defin{contexte de typage} est un ensemble fini $\Gamma$ de typages
$x_i:\sigma_i$, où les $x_i$ sont des variables de terme distinctes.
On l'écrit généralement
\[
\Gamma=x_1:\sigma_1,\ldots,x_k:\sigma_k.
\]

Un \defin{jugement de typage} est une expression
\[
\Gamma\vdash E:\tau,
\]
signifiant que $E$ est un terme de type $\tau$ dans le contexte
$\Gamma$. Lorsque $\Gamma$ est vide, on écrit simplement
$\vdash E:\tau$.
\end{defn}

\thingy Les règles de typage du $\lambda$CST sont les suivantes :
\[
\inferrule*[Left={Var}]
 { }
 {\Gamma,x:\sigma\vdash x:\sigma}
\]
\[
\inferrule*[Left={App},sep=4em]
 {\Gamma\vdash P:\sigma\to\tau\\
  \Gamma\vdash Q:\sigma}
 {\Gamma\vdash PQ:\tau}
\]
\[
\inferrule*[Left={Abs}]
 {\Gamma,v:\sigma\vdash E:\tau}
 {\Gamma\vdash\lambda(v:\sigma).E:\sigma\to\tau}.
\]
Chaque barre horizontale signifie que le jugement placé en dessous
découle, par la règle indiquée, des jugements placés au-dessus.
L'ensemble des jugements de typage est le plus petit ensemble fermé
par ces règles.

Une \defin{dérivation} d'un jugement est un arbre d'instances de ces
règles dont la racine est le jugement considéré.

\thingy Voici quelques exemples de jugements :
\[
f:\alpha\to\beta,\;x:\alpha\vdash fx:\beta,
\]
\[
f:\beta\to\alpha\to\gamma,\;x:\alpha,\;y:\beta
  \vdash fyx:\gamma,
\]
\[
f:\beta\to\alpha\to\gamma,\;x:\alpha
  \vdash\lambda(y:\beta).fyx:\beta\to\gamma,
\]
\[
x:\alpha\vdash
  \lambda(f:\alpha\to\beta).fx:(\alpha\to\beta)\to\beta,
\]
et
\[
\vdash
\lambda(x:\alpha).\lambda(f:\alpha\to\beta).fx:
\alpha\to(\alpha\to\beta)\to\beta.
\]

À titre d'exemple plus développé, on obtient la dérivation suivante :
{\footnotesize
\[
\inferrule*[Left={Abs}]{
 \inferrule*[Left={Abs}]{
  \inferrule*[Left={Abs}]{
   \inferrule*[Left={App}]{
    \inferrule*[Left={App}]{
     \inferrule*[Left={Var}]{ }
      {f:\beta\to\alpha\to\gamma,x:\alpha,y:\beta
       \vdash f:\beta\to\alpha\to\gamma}
     \\
     \inferrule*[Right={Var}]{ }
      {f:\beta\to\alpha\to\gamma,x:\alpha,y:\beta
       \vdash y:\beta}
    }{
     f:\beta\to\alpha\to\gamma,x:\alpha,y:\beta
     \vdash fy:\alpha\to\gamma
    }
    \\
    \inferrule*[Right={Var}]{ }
     {f:\beta\to\alpha\to\gamma,x:\alpha,y:\beta
      \vdash x:\alpha}
   }{
    f:\beta\to\alpha\to\gamma,x:\alpha,y:\beta
    \vdash fyx:\gamma
   }
  }{
   f:\beta\to\alpha\to\gamma,x:\alpha
   \vdash\lambda(y:\beta).fyx:\beta\to\gamma
  }
 }{
  f:\beta\to\alpha\to\gamma
  \vdash\lambda(x:\alpha).\lambda(y:\beta).fyx:
  \alpha\to\beta\to\gamma
 }
}{
 \vdash
 \lambda(f:\beta\to\alpha\to\gamma).
 \lambda(x:\alpha).\lambda(y:\beta).fyx:
 (\beta\to\alpha\to\gamma)\to(\alpha\to\beta\to\gamma)
}
\]
}

\subsection{Propriétés élémentaires du typage}

\begin{prop}[Affaiblissement]
Si $\Gamma\subseteq\Gamma'$ sont deux contextes et si
$\Gamma\vdash M:\sigma$, alors $\Gamma'\vdash M:\sigma$.
\end{prop}

\begin{proof}
On raisonne par induction sur la dérivation de
$\Gamma\vdash M:\sigma$. Dans le cas d'une variable, le typage utilisé
appartient à $\Gamma$, donc également à $\Gamma'$. Les règles
d'application et d'abstraction se transportent immédiatement en
appliquant l'hypothèse d'induction à leurs prémisses.
\end{proof}

\begin{prop}[Variables libres]
Si
\[
x_1:\sigma_1,\ldots,x_k:\sigma_k\vdash E:\tau,
\]
alors toute variable libre de $E$ appartient à
$\{x_1,\ldots,x_k\}$.
\end{prop}

\begin{proof}
La propriété se vérifie par induction sur la dérivation. La règle
variable ne peut introduire qu'une variable présente dans le contexte.
L'application réunit les variables libres de ses deux sous-termes.
Enfin, l'abstraction lie précisément la variable ajoutée au contexte.
\end{proof}

En conséquence, les typages du contexte qui ne concernent aucune
variable libre de $E$ peuvent être supprimés. Le typage est également
invariant par renommage cohérent des variables libres.

\thingy La dérivation du typage d'un préterme se construit
systématiquement à partir de sa syntaxe :
\begin{itemize}
\item une variable doit être présente dans le contexte ;
\item pour une application $PQ$, les types obtenus pour $P$ et $Q$
  doivent être respectivement $\sigma\to\tau$ et $\sigma$ ;
\item pour une abstraction $\lambda(v:\sigma).E$, on type $E$ dans le
  contexte enrichi par $v:\sigma$.
\end{itemize}
Il n'y a donc aucun choix substantiel à effectuer.

\begin{prop}[Unicité du type]
Pour un contexte $\Gamma$ et un préterme annoté $M$ donnés, il existe
au plus un type $\sigma$ tel que $\Gamma\vdash M:\sigma$.
\end{prop}

\begin{proof}
On raisonne par induction sur la structure de $M$. Le type d'une
variable est imposé par le contexte. Celui d'une abstraction est imposé
par l'annotation de sa variable et par le type, unique par hypothèse
d'induction, de son corps. Enfin, dans une application, le type de la
fonction et celui de l'argument imposent de manière unique le type du
résultat.
\end{proof}

\subsection{Substitution, réduction et normalisation}

\begin{lem}[Substitution]\label{substitution-lemma}
Si
\[
\Gamma,v:\sigma\vdash E:\tau
\qquad\text{et}\qquad
\Gamma\vdash T:\sigma,
\]
alors
\[
\Gamma\vdash E[v\backslash T]:\tau,
\]
où $E[v\backslash T]$ désigne la substitution correcte de $v$ par $T$
dans $E$.
\end{lem}

\begin{proof}
On raisonne par induction sur la dérivation de
$\Gamma,v:\sigma\vdash E:\tau$. Les règles d'application et
d'abstraction se transportent en appliquant l'hypothèse d'induction aux
prémisses, après avoir renommé si nécessaire les variables liées pour
éviter toute capture.

Dans le cas de la règle variable, si la variable considérée est $v$,
on remplace la feuille correspondante par la dérivation de
$\Gamma\vdash T:\sigma$. Si elle est différente de $v$, la même règle
variable reste applicable dans $\Gamma$. Ceci construit une dérivation
de $\Gamma\vdash E[v\backslash T]:\tau$.
\end{proof}

\thingy La $\beta$-réduction remplace un redex
\[
(\lambda(v:\sigma).E)T
\]
par son réduit
\[
E[v\backslash T].
\]
On note $X\rightarrow_\beta X'$ pour une étape de
$\beta$-réduction et $X\twoheadrightarrow_\beta X'$ pour une suite de
telles étapes.

\begin{prop}[Préservation du typage]
Si $\Gamma\vdash X:\tau$ et
$X\twoheadrightarrow_\beta X'$, alors
$\Gamma\vdash X':\tau$.
\end{prop}

\begin{proof}
Il suffit de considérer une étape de réduction. Lorsqu'un redex
$(\lambda(v:\sigma).E)T$ est bien typé, on dispose de
\[
\Gamma,v:\sigma\vdash E:\tau
\qquad\text{et}\qquad
\Gamma\vdash T:\sigma.
\]
Le lemme de substitution donne alors
$\Gamma\vdash E[v\backslash T]:\tau$. La réduction à l'intérieur d'un
contexte se traite par induction sur ce contexte. Une suite de
réductions se traite ensuite par induction sur sa longueur.
\end{proof}

\begin{thm}[Turing, Tait, Girard]
Le $\lambda$-calcul simplement typé est \defin{fortement normalisant} :
toute suite de $\beta$-réductions issue d'un terme bien typé est finie.
\end{thm}

\begin{proof}
Les diapositives indiquent la méthode dite de calculabilité. On définit
par induction sur le type la notion de terme \emph{fortement
calculable} :
\begin{itemize}
\item un terme de type atomique est fortement calculable lorsqu'il est
  fortement normalisable ;
\item un terme $P$ de type $\sigma\to\tau$ est fortement calculable
  lorsque, pour tout terme $Q$ de type $\sigma$ fortement calculable,
  le terme $PQ$ est fortement calculable.
\end{itemize}
On montre d'abord, par induction sur le type, que tout terme fortement
calculable est fortement normalisable. On montre ensuite, par induction
sur une dérivation de typage, que la substitution de termes fortement
calculables aux variables libres d'un terme bien typé produit un terme
fortement calculable. En substituant les variables elles-mêmes, on
obtient que tout terme bien typé est fortement calculable, donc
fortement normalisable.

\review{La preuve précédente développe seulement l'esquisse donnée dans
les diapositives. Une version finale devrait préciser les propriétés de
stabilité des termes fortement calculables utilisées dans le cas de
l'abstraction.}
\end{proof}

\subsection{Reconstruction des termes et problèmes algorithmiques}

\thingy La syntaxe du terme détermine directement la forme de sa
dérivation de typage. Réciproquement, si l'on connaît un arbre de
dérivation dans lequel les types et les règles sont indiqués, mais où
les termes ont été effacés, les termes peuvent être reconstruits en
suivant l'arbre :
\begin{itemize}
\item une règle variable détermine la variable utilisée ;
\item une règle d'application reconstruit l'application des deux termes
  obtenus dans les prémisses ;
\item une règle d'abstraction reconstruit la variable liée et le corps.
\end{itemize}
Il faut seulement distinguer correctement les variables lorsque
plusieurs hypothèses du contexte ont le même type.

\thingy Les problèmes suivants ont des difficultés différentes :
\begin{itemize}
\item étant donné un préterme annoté, calculer son type ou vérifier un
  jugement est facile et décidable ;
\item étant donné un arbre de dérivation typé dont les termes ont été
  effacés, reconstruire le terme est également facile ;
\item étant donné un terme non annoté, déterminer s'il admet des
  annotations de type demande un véritable algorithme d'inférence ;
\item étant donné un type, déterminer s'il est habité correspond, par
  Curry--Howard, à la décidabilité du calcul propositionnel
  intuitionniste.
\end{itemize}

\section{Le calcul propositionnel intuitionniste}

\subsection{Formules et déduction naturelle}

\thingy Le \defin{calcul propositionnel} est une logique dont les
objets de base sont des propositions, sans variables d'individus ni
quantificateurs $\forall$ ou $\exists$. Ses connecteurs sont
\[
\Rightarrow,\qquad \land,\qquad \lor,\qquad \top,\qquad \bot.
\]
La négation est une abréviation :
\[
\neg P:=P\Rightarrow\bot.
\]

Une \defin{formule propositionnelle} est définie inductivement comme
une variable propositionnelle, ou comme l'application d'un connecteur
aux formules appropriées. On adopte les priorités usuelles : $\neg$
est plus prioritaire que $\land$, qui est plus prioritaire que $\lor$,
qui est plus prioritaire que $\Rightarrow$. L'implication associe à
droite.

Un \defin{séquent}
\[
P_1,\ldots,P_r\vdash Q
\]
signifie que $Q$ est démontrable sous les hypothèses
$P_1,\ldots,P_r$.

\thingy La logique classique admet le tiers exclu, sous l'une de ses
formes équivalentes. La logique intuitionniste étudiée ici ne l'admet
pas. Il ne faut pas l'interpréter comme une logique possédant une
troisième valeur de vérité : elle impose plutôt une notion plus
constructive de preuve.

Chaque connecteur possède des règles d'\defin{introduction}, permettant
de le démontrer, et des règles d'\defin{élimination}, permettant de
l'utiliser.

Les règles de la déduction naturelle intuitionniste sont :
\[
\inferrule*[Left={Ax}]
 { }
 {\Gamma,P\vdash P}
\]
\[
\inferrule*[Left={$\Rightarrow$Int}]
 {\Gamma,P\vdash Q}
 {\Gamma\vdash P\Rightarrow Q}
\qquad
\inferrule*[Left={$\Rightarrow$Élim}]
 {\Gamma\vdash P\Rightarrow Q\\
  \Gamma\vdash P}
 {\Gamma\vdash Q}
\]
\[
\inferrule*[Left={$\land$Int}]
 {\Gamma\vdash Q_1\\
  \Gamma\vdash Q_2}
 {\Gamma\vdash Q_1\land Q_2}
\]
\[
\inferrule*[Left={$\land$Élim$_1$}]
 {\Gamma\vdash Q_1\land Q_2}
 {\Gamma\vdash Q_1}
\qquad
\inferrule*[Left={$\land$Élim$_2$}]
 {\Gamma\vdash Q_1\land Q_2}
 {\Gamma\vdash Q_2}
\]
\[
\inferrule*[Left={$\lor$Int$_1$}]
 {\Gamma\vdash Q_1}
 {\Gamma\vdash Q_1\lor Q_2}
\qquad
\inferrule*[Left={$\lor$Int$_2$}]
 {\Gamma\vdash Q_2}
 {\Gamma\vdash Q_1\lor Q_2}
\]
\[
\inferrule*[Left={$\lor$Élim}]
 {\Gamma\vdash P_1\lor P_2\\
  \Gamma,P_1\vdash Q\\
  \Gamma,P_2\vdash Q}
 {\Gamma\vdash Q}
\]
\[
\inferrule*[Left={$\top$Int}]
 { }
 {\Gamma\vdash\top}
\qquad
\inferrule*[Left={$\bot$Élim}]
 {\Gamma\vdash\bot}
 {\Gamma\vdash Q}.
\]

Dans ces règles, les hypothèses mises en évidence dans une prémisse
d'une règle d'introduction de $\Rightarrow$ ou d'élimination de $\lor$
sont déchargées dans la conclusion.

\thingy Par exemple, la commutativité de la conjonction se démontre par
\[
\inferrule*[Left={$\Rightarrow$Int}]{
 \inferrule*[Left={$\land$Int}]{
  \inferrule*[Left={$\land$Élim$_2$}]{
   \inferrule*[Left={Ax}]{ }{A\land B\vdash A\land B}
  }{A\land B\vdash B}
  \\
  \inferrule*[Right={$\land$Élim$_1$}]{
   \inferrule*[Right={Ax}]{ }{A\land B\vdash A\land B}
  }{A\land B\vdash A}
 }{A\land B\vdash B\land A}
}{\vdash A\land B\Rightarrow B\land A}.
\]

La commutativité de la disjonction se démontre en distinguant les deux
cas fournis par une preuve de $A\lor B$ :
\[
\inferrule*[Left={$\Rightarrow$Int}]{
 \inferrule*[Left={$\lor$Élim}]{
  \inferrule*[Left={Ax}]{ }{A\lor B\vdash A\lor B}
  \\
  \inferrule*[Right={$\lor$Int$_2$}]{
   \inferrule*[Right={Ax}]{ }{A\lor B,A\vdash A}
  }{A\lor B,A\vdash B\lor A}
  \\
  \inferrule*[Right={$\lor$Int$_1$}]{
   \inferrule*[Right={Ax}]{ }{A\lor B,B\vdash B}
  }{A\lor B,B\vdash B\lor A}
 }{A\lor B\vdash B\lor A}
}{\vdash A\lor B\Rightarrow B\lor A}.
\]

\subsection{Théorèmes et non-théorèmes}

\thingy Une formule $P$ telle que $\vdash P$ est dite
\defin{valable}, ou encore un \defin{théorème} du calcul
propositionnel intuitionniste.

Parmi les théorèmes importants figurent
\[
A\Rightarrow A,
\qquad
A\Rightarrow B\Rightarrow A,
\]
\[
(A\Rightarrow B\Rightarrow C)
\Rightarrow(A\Rightarrow B)\Rightarrow A\Rightarrow C,
\]
ainsi que les lois usuelles d'associativité, de commutativité et de
distributivité de $\land$ et $\lor$, sous leurs formes
intuitionnistement valables.

En revanche, les implications marquées dans les diapositives comme
classiques ne sont pas nécessairement démontrables
intuitionnistement. Parmi les non-théorèmes figurent :
\[
((A\Rightarrow B)\Rightarrow A)\Rightarrow A,
\]
\[
A\lor\neg A,
\qquad
\neg\neg A\Rightarrow A,
\]
\[
\neg(A\land B)\Rightarrow\neg A\lor\neg B,
\]
\[
(A\Rightarrow B)\Rightarrow\neg A\lor B,
\]
et
\[
(\neg B\Rightarrow\neg A)\Rightarrow(A\Rightarrow B).
\]
Ces formules sont valables en logique classique, mais pas dans la
logique intuitionniste.

\thingy Il faut distinguer une preuve de la négation d'une preuve par
l'absurde.

Pour démontrer $\neg A$, on suppose $A$ et on en déduit $\bot$. Cette
démonstration est intuitionnistement valable, puisqu'elle construit une
preuve de $A\Rightarrow\bot$.

Une preuve classique par l'absurde suppose au contraire $\neg A$,
obtient une contradiction, puis conclut $A$. Intuitionnistement, cet
argument démontre seulement $\neg\neg A$, et non nécessairement $A$.
Ainsi, $\neg A$ peut se lire « $A$ conduit à l'absurde », tandis que
$\neg\neg A$ signifie seulement que la négation de $A$ conduit à
l'absurde.

\subsection{Logique classique}

\thingy La logique classique peut être obtenue en remplaçant la règle
d'élimination de $\bot$ par une règle de raisonnement par l'absurde :
\[
\inferrule*[Left={Absurde}]
 {\Gamma,\neg Q\vdash\bot}
 {\Gamma\vdash Q}.
\]
Elle possède donc davantage de théorèmes que la logique
intuitionniste.

La sémantique booléenne fournit une caractérisation de la logique
classique : une formule $P$ est un théorème classique si et seulement
si toute affectation de $\bot$ ou $\top$ à ses variables lui donne la
valeur $\top$.

\begin{thm}[Correction et complétude classiques]
Les tables de vérité booléennes sont correctes et complètes pour le
calcul propositionnel classique.
\end{thm}

La correction se démontre par induction sur les dérivations. Pour la
complétude, on peut raisonner successivement sur les valeurs possibles
des variables et utiliser le tiers exclu afin de couvrir toutes les
lignes de la table de vérité.

\subsection{Interprétation de Brouwer--Heyting--Kolmogorov}

\thingy L'interprétation de Brouwer--Heyting--Kolmogorov donne une
lecture constructive des connecteurs :
\begin{itemize}
\item un témoignage de $P\land Q$ est un témoignage de $P$ accompagné
  d'un témoignage de $Q$ ;
\item un témoignage de $P\lor Q$ est un témoignage de l'un des deux,
  accompagné de l'indication du cas choisi ;
\item un témoignage de $P\Rightarrow Q$ est un moyen de transformer
  tout témoignage de $P$ en un témoignage de $Q$ ;
\item un témoignage de $\top$ est trivial ;
\item il n'existe aucun témoignage de $\bot$.
\end{itemize}
Cette interprétation explique directement les règles de déduction
naturelle.

\subsection{Correspondance de Curry--Howard}

\thingy Pour le connecteur $\Rightarrow$, les règles du $\lambda$CST et
celles de la déduction naturelle sont exactement les mêmes, au
changement de vocabulaire près :
\[
\begin{array}{c|c}
\text{Typage du $\lambda$CST}
&
\text{Calcul propositionnel intuitionniste}
\\[1ex]\hline\\[-2ex]
\inferrule*[Left={Var}]
 { }
 {\Gamma,x:\sigma\vdash x:\sigma}
&
\inferrule*[Left={Ax}]
 { }
 {\Gamma,P\vdash P}
\\[3ex]
\inferrule*[Left={App}]
 {\Gamma\vdash f:\sigma\to\tau\\
  \Gamma\vdash x:\sigma}
 {\Gamma\vdash fx:\tau}
&
\inferrule*[Left={$\Rightarrow$Élim}]
 {\Gamma\vdash P\Rightarrow Q\\
  \Gamma\vdash P}
 {\Gamma\vdash Q}
\\[3ex]
\inferrule*[Left={Abs}]
 {\Gamma,v:\sigma\vdash t:\tau}
 {\Gamma\vdash\lambda(v:\sigma).t:\sigma\to\tau}
&
\inferrule*[Left={$\Rightarrow$Int}]
 {\Gamma,P\vdash Q}
 {\Gamma\vdash P\Rightarrow Q}
\end{array}
\]
On peut donc identifier les termes du $\lambda$CST aux preuves en
déduction naturelle restreinte à l'implication.

Par exemple, le terme
\[
\lambda(f:\beta\to\alpha\to\gamma).
\lambda(x:\alpha).
\lambda(y:\beta).fyx
\]
correspond à une preuve de
\[
(B\Rightarrow A\Rightarrow C)
\Rightarrow(A\Rightarrow B\Rightarrow C).
\]

\thingy Pour traiter la conjonction, on enrichit le $\lambda$CST de
types produits. Les règles sont
\[
\inferrule*[Left={Pair}]
 {\Gamma\vdash t_1:\tau_1\\
  \Gamma\vdash t_2:\tau_2}
 {\Gamma\vdash\langle t_1,t_2\rangle:\tau_1\times\tau_2}
\]
et
\[
\inferrule*[Left={Proj$_i$}]
 {\Gamma\vdash t:\tau_1\times\tau_2}
 {\Gamma\vdash\pi_i t:\tau_i}
 \qquad(i\in\{1,2\}).
\]
Elles correspondent exactement aux règles d'introduction et
d'élimination de $\land$. Les réductions associées sont
\[
\pi_i\langle t_1,t_2\rangle\rightsquigarrow t_i.
\]
La preuve de $A\land B\Rightarrow B\land A$ correspond ainsi au terme
\[
\lambda(u:\alpha\times\beta).
\langle\pi_2u,\pi_1u\rangle.
\]

Pour traiter la disjonction, on ajoute des types sommes, avec deux
injections
\[
\iota_i^{(\tau_1,\tau_2)}:\tau_i\to\tau_1+\tau_2
\qquad(i\in\{1,2\}),
\]
et une construction par analyse de cas :
\[
\inferrule*[Left={Match}]
 {\Gamma\vdash r:\sigma_1+\sigma_2\\
  \Gamma,v_1:\sigma_1\vdash t_1:\tau\\
  \Gamma,v_2:\sigma_2\vdash t_2:\tau}
 {\Gamma\vdash
  (\texttt{match~}r\texttt{~with~}
   \iota_1v_1\mapsto t_1,\;
   \iota_2v_2\mapsto t_2):\tau}.
\]
Les réductions associées sont
\[
(\texttt{match~}\iota_i s\texttt{~with~}
 \iota_1v_1\mapsto t_1,\;
 \iota_2v_2\mapsto t_2)
\rightsquigarrow t_i[v_i\backslash s].
\]

Enfin, $\top$ correspond à un type unité $1$ possédant une unique
valeur $\bullet$, tandis que $\bot$ correspond à un type vide $0$.
Une élimination du type vide fournit un terme de n'importe quel type :
\[
\inferrule*[Left={Void}]
 {\Gamma\vdash r:0}
 {\Gamma\vdash\texttt{exfalso}^{(\tau)}r:\tau}.
\]

Ces extensions restent fortement normalisantes.

\thingy La correspondance de Curry--Howard peut alors se résumer ainsi :
\begin{center}
\begin{tabular}{c|c}
Types & Propositions\\ \hline
$\sigma\to\tau$ & $P\Rightarrow Q$\\
$\sigma\times\tau$ & $P\land Q$\\
$\sigma+\tau$ & $P\lor Q$\\
$1$ & $\top$\\
$0$ & $\bot$
\end{tabular}
\end{center}
Les termes correspondent aux preuves, l'abstraction à l'ouverture
d'une hypothèse, l'application au \textit{modus ponens}, et les
variables liées aux hypothèses déchargées.

Une même proposition peut avoir plusieurs preuves fonctionnellement
différentes. Par exemple, les entiers de Church
\[
\overline{0}_\alpha
 :=\lambda(f:\alpha\to\alpha).\lambda(x:\alpha).x,
\]
\[
\overline{1}_\alpha
 :=\lambda(f:\alpha\to\alpha).\lambda(x:\alpha).fx,
\]
\[
\overline{2}_\alpha
 :=\lambda(f:\alpha\to\alpha).\lambda(x:\alpha).f(fx)
\]
ont tous le type
\[
(\alpha\to\alpha)\to(\alpha\to\alpha)
\]
et correspondent donc à différentes preuves du même énoncé
$(A\Rightarrow A)\Rightarrow(A\Rightarrow A)$.

\subsection{Calcul des séquents et élimination des coupures}

\thingy Le \defin{calcul des séquents} est une autre présentation de la
même logique propositionnelle intuitionniste. Il distingue, pour
chaque connecteur, une règle droite qui l'introduit dans la conclusion
et une règle gauche qui l'introduit dans les hypothèses. On conserve
notamment l'axiome
\[
\inferrule*[Left={Ax}]{ }{\Gamma,Q\vdash Q},
\]
ainsi que les règles structurelles
\[
\inferrule*[Left={Contr}]
 {\Gamma,P,P\vdash Q}
 {\Gamma,P\vdash Q}
\qquad
\inferrule*[Left={Weak}]
 {\Gamma\vdash Q}
 {\Gamma,P\vdash Q}.
\]

La règle de \defin{coupure} est
\[
\inferrule*[Left={Cut}]
 {\Gamma\vdash P\\
  \Gamma,P\vdash Q}
 {\Gamma\vdash Q}.
\]
Elle exprime la composition de deux démonstrations, $P$ étant la
formule coupée.

\begin{thm}[Gentzen]
La règle de coupure est éliminable : tout séquent démontrable avec la
règle de coupure est démontrable sans elle.
\end{thm}

\begin{proof}
La preuve transforme localement une démonstration contenant une
coupure. On effectue une double récurrence : d'abord sur le degré de la
formule coupée, puis, à degré fixé, sur la somme des hauteurs des deux
branches de la coupure.

Lorsque la dernière règle de l'une des branches ne concerne pas la
formule coupée, on fait remonter la coupure au-dessus de cette règle.
Lorsqu'une branche se termine par l'axiome portant sur la formule
coupée, la coupure disparaît. Enfin, lorsque les deux dernières règles
introduisent le connecteur principal de la formule coupée, on remplace
la coupure par une ou plusieurs coupures portant sur des sous-formules
strictes. Le degré diminue alors, ce qui permet d'appliquer l'hypothèse
de récurrence.

\review{La preuve complète demande l'examen systématique de toutes les
paires de règles principales. Les diapositives n'en donnent que la
méthode et quelques exemples de transformations locales.}
\end{proof}

\thingy L'élimination des coupures a plusieurs conséquences.

Elle implique la \defin{propriété de la disjonction} :
\[
\text{si }\vdash Q_1\lor Q_2,\quad
\text{alors }\vdash Q_1\text{ ou }\vdash Q_2.
\]
En effet, une démonstration sans coupure d'une disjonction sans
hypothèse doit se terminer par l'une des deux règles d'introduction de
$\lor$.

Elle implique également la \defin{propriété de la sous-formule} : une
preuve sans coupure de
\[
P_1,\ldots,P_r\vdash Q
\]
ne fait intervenir que des sous-formules des $P_i$ et de $Q$.

On en déduit un algorithme de décision du calcul propositionnel
intuitionniste : il suffit d'énumérer les séquents construits avec les
sous-formules pertinentes et de fermer cet ensemble par les règles
autres que la coupure.

\section{Le call/cc et la logique classique}

\subsection{Continuations}

\thingy Dans un langage de programmation, la \defin{continuation} d'un
appel de fonction est l'état du calcul qui attend que cette fonction
renvoie une valeur. On peut la voir comme une copie de la pile d'appels
au moment de l'appel.

Certains langages permettent de capturer cette continuation, de la
stocker comme une valeur et de l'invoquer explicitement. Invoquer une
continuation avec une valeur $v$ restaure la pile capturée et fait
comme si l'appel concerné avait renvoyé $v$, même si ce calcul avait
déjà terminé. Une continuation est ainsi comparable à un
\texttt{goto} qui restaure en même temps la pile d'appels.

Les continuations peuvent remplacer certains mécanismes d'exception,
mais aussi reprendre un calcul interrompu. Elles permettent notamment
d'implémenter des générateurs et certaines formes de multitâche
coopératif. Leur coût d'implémentation porte essentiellement sur la
gestion de la pile.

\thingy La fonction \defin{call/cc}, abréviation de
\textit{call-with-current-continuation}, prend en argument une fonction
$g$, capture sa propre continuation et passe celle-ci à $g$.

Une continuation ne retourne jamais au point où elle est invoquée.
Elle peut donc être typée comme $\alpha\to\bot$, ou plus généralement
comme $\alpha\to\beta$ pour un type de retour arbitraire $\beta$.
La fonction $g$ attend une continuation de type $\alpha\to\beta$ et
renvoie une valeur de type $\alpha$. On obtient ainsi
\[
\texttt{call/cc}:
((\alpha\to\beta)\to\alpha)\to\alpha.
\]
Ce type correspond à la loi de Peirce
\[
((A\Rightarrow B)\Rightarrow A)\Rightarrow A,
\]
qui est une formule caractéristique de la logique classique.

La présence de \texttt{call/cc} transforme donc la logique du typage :
la correspondance de Curry--Howard ne donne plus seulement une logique
intuitionniste, mais une logique classique. Cette interprétation peut
être formalisée par le $\lambda\mu$-calcul.

\thingy Le tiers exclu
\[
A\lor\neg A
\]
peut alors être implémenté en capturant une continuation. L'idée est de
renvoyer provisoirement le cas $\neg A$. Si cette prétendue négation
est ensuite appliquée à une valeur de type $A$, la continuation
capturée est invoquée afin de revenir en arrière et de renvoyer cette
fois le cas $A$.

Cette possibilité montre pourquoi la logique classique est moins
constructive : une valeur somme annoncée comme appartenant à l'un des
deux cas n'est plus stable, puisqu'une continuation peut faire revenir
le calcul dans le passé et modifier le choix.

\subsection{Continuation Passing Style}

\thingy Même lorsqu'un langage ne fournit pas directement de
continuations de première classe, on peut réécrire les programmes en
\defin{Continuation Passing Style}, ou CPS. Au lieu de renvoyer une
valeur $v$, une fonction reçoit en argument une continuation $k$ et
appelle $k$ sur $v$. Tous les appels deviennent ainsi terminaux.

Fixons un type $Z$ de retour ultime et posons
\[
\mathop{\sim}P:=(P\Rightarrow Z),
\qquad
\mathop{\sim}\mathop{\sim}P
  :=((P\Rightarrow Z)\Rightarrow Z).
\]
Une valeur de type $P$ passée par continuation a donc pour type
$\mathop{\sim}\mathop{\sim}P$.

Pour les termes du $\lambda$-calcul non typé, la transformation
essentielle est donnée par
\[
v^{\mathrm{CPS}}:=\lambda k.kv,
\]
\[
(\lambda v.t)^{\mathrm{CPS}}
 :=\lambda k.\,k(\lambda v.t^{\mathrm{CPS}}),
\]
et
\[
(fx)^{\mathrm{CPS}}
 :=\lambda k.\,
 f^{\mathrm{CPS}}
 \bigl(\lambda f_0.\,
 x^{\mathrm{CPS}}
 \bigl(\lambda x_0.\,f_0x_0k\bigr)\bigr).
\]
Cette dernière formule évalue d'abord $f$, puis $x$, puis applique la
fonction obtenue à la valeur obtenue et à la continuation du calcul
complet.

\thingy Pour typer systématiquement cette transformation, on définit
d'abord $P^\diamond$ par
\[
A^\diamond=A
\]
pour une variable de type $A$, et
\[
(P\Rightarrow Q)^\diamond
 :=P^\diamond\Rightarrow
   \mathop{\sim}\mathop{\sim}Q^\diamond.
\]
Les produits, sommes, $\top$ et $\bot$ sont transformés composante par
composante. On pose ensuite
\[
P^{\mathrm{CPS}}
 :=\mathop{\sim}\mathop{\sim}P^\diamond.
\]
La transformation des termes est définie de manière à garantir
\[
\text{si }\vdash t:P,
\qquad
\text{alors }\vdash t^{\mathrm{CPS}}:P^{\mathrm{CPS}}.
\]

Lorsque $Z=\bot$, cette transformation correspond à une traduction de
la logique classique dans la logique intuitionniste par double
négation. Elle donne notamment
\[
\mathsf{CPC}\vdash P
\quad\Longleftrightarrow\quad
\mathsf{IPC}\vdash P^{\mathrm{CPS}}.
\]
\review{Les diapositives indiquent plusieurs variantes de la traduction
par double négation, notamment celles de Gödel--Gentzen et de Glivenko.
Il faudra décider dans une version ultérieure jusqu'où développer leur
comparaison.}

\section{Sémantique du calcul propositionnel intuitionniste}

\subsection{Sémantiques, correction et complétude}

\thingy Une logique $\mathsf{L}$ est définie syntaxiquement par des
formules, des axiomes et des règles de déduction. On note
$\mathsf{L}\vdash\varphi$ lorsque $\varphi$ est démontrable dans cette
logique.

Une \defin{sémantique} est une classe de modèles donnant un sens aux
symboles de la logique. On note
\[
\mathcal{M}\vDash\varphi
\]
lorsque la formule $\varphi$ est vraie dans le modèle $\mathcal{M}$.

Une sémantique est \defin{correcte} lorsque tout théorème est vrai dans
tout modèle :
\[
\mathsf{L}\vdash\varphi
\quad\Longrightarrow\quad
\mathsf{S}\vDash\varphi.
\]
Elle est \defin{complète} lorsque toute formule vraie dans tous les
modèles est démontrable :
\[
\mathsf{S}\vDash\varphi
\quad\Longrightarrow\quad
\mathsf{L}\vdash\varphi.
\]

La sémantique booléenne est correcte et complète pour la logique
classique, mais elle n'est que correcte pour la logique
intuitionniste. Il faut donc chercher des modèles plus riches pour
caractériser exactement cette dernière.

\subsection{Cadres de Kripke}

\begin{defn}
Un \defin{cadre de Kripke} est un ensemble partiellement ordonné
$(W,\leq)$. Ses éléments sont appelés des \defin{mondes}. Si
$w\leq w'$, on dit que $w'$ est accessible depuis $w$.

Une affectation de vérité sur ce cadre associe à chaque variable
propositionnelle une fonction
\[
p:W\to\{0,1\}
\]
croissante : si $p(w)=1$ et $w\leq w'$, alors $p(w')=1$. Cette
propriété est appelée permanence.

Un modèle de Kripke est un cadre muni d'une telle affectation pour
chaque variable propositionnelle.
\end{defn}

\thingy Les connecteurs sont interprétés dans un monde $w$ par
\begin{itemize}
\item $(q_1\mathbin{\dottedland}q_2)(w)=1$ si et seulement si
  $q_1(w)=q_2(w)=1$ ;
\item $(q_1\mathbin{\dottedlor}q_2)(w)=1$ si et seulement si
  $q_1(w)=1$ ou $q_2(w)=1$ ;
\item $(q_1\mathbin{\dottedlimp}q_2)(w)=1$ si et seulement si, pour
  tout $w'\geq w$, $q_1(w')=1$ implique $q_2(w')=1$ ;
\item $\dottedtop(w)=1$ et $\dottedbot(w)=0$.
\end{itemize}
En particulier,
\[
\dot{\neg}p
 :=p\mathbin{\dottedlimp}\dottedbot
\]
vaut dans $w$ lorsque $p$ est faux dans tout monde accessible depuis
$w$.

Ainsi, $A\Rightarrow B$ ne signifie pas simplement que $A$ est faux ou
que $B$ est vrai dans le monde actuel. Il signifie que, dans tout
monde futur où $A$ deviendra vrai, $B$ sera également vrai.

Une formule est valide dans un modèle lorsqu'elle vaut dans tous ses
mondes, et valide pour la sémantique de Kripke lorsqu'elle vaut dans
tout modèle de Kripke.

\begin{thm}[Kripke]
La sémantique de Kripke est correcte et complète pour le calcul
propositionnel intuitionniste.
\end{thm}

\review{La preuve du théorème de Kripke n'est pas donnée dans les
diapositives et n'est donc pas développée ici.}

\thingy La complétude exige de considérer tous les cadres. Un cadre
réduit à un singleton redonne la sémantique booléenne. Un cadre
totalement ordonné valide
\[
(A\Rightarrow B)\lor(B\Rightarrow A),
\]
qui n'est pourtant pas démontrable intuitionnistement en général.

Le cadre à deux mondes $u\leq v$ donne une logique à trois valeurs
$\{0,\mathsf{q},1\}$, où $\mathsf{q}$ signifie qu'une proposition est
fausse dans $u$ mais vraie dans $v$.

\subsection{Sémantique topologique}

\thingy Fixons un espace topologique $X$ et notons $\mathcal{O}(X)$ son
ensemble d'ouverts. On interprète les connecteurs sur les ouverts par
\[
U\mathbin{\dottedland}V:=U\cap V,
\qquad
U\mathbin{\dottedlor}V:=U\cup V,
\]
\[
\dottedtop:=X,
\qquad
\dottedbot:=\varnothing,
\]
et
\[
U\mathbin{\dottedlimp}V
 :=\operatorname{int}\bigl((X\setminus U)\cup V\bigr).
\]
Autrement dit, $U\mathbin{\dottedlimp}V$ est le plus grand ouvert $W$
tel que
\[
U\cap W\subseteq V.
\]
On obtient notamment
\[
\dot{\neg}U=\operatorname{int}(X\setminus U)
\]
et
\[
\dot{\neg}\dot{\neg}U
 =\operatorname{int}(\operatorname{adh}(U)).
\]

Une formule est interprétée en substituant des ouverts à ses variables.
Elle est valide si son interprétation est toujours l'ouvert $X$.

\begin{thm}[Tarski]
La sémantique des ouverts est correcte et complète pour le calcul
propositionnel intuitionniste. Pour la complétude, il suffit même de
considérer les espaces $\mathbb{R}^n$, pour $n\geq1$.
\end{thm}

Cette sémantique traduit l'intuition suivant laquelle la vérité est
locale : si une propriété est vraie en un point, elle doit rester vraie
dans un voisinage de ce point.

\thingy À titre d'exemple, dans $X=\mathbb{R}$, soit
$U={]0,1[}$. Alors
\[
\dot{\neg}U
 ={]-\infty,0[}\cup{]1,+\infty[},
\qquad
\dot{\neg}\dot{\neg}U=U.
\]
On vérifie également que
\[
(\dot{\neg}\dot{\neg}U\mathbin{\dottedlimp}U)=\mathbb{R}.
\]
En revanche,
\[
U\mathbin{\dottedlor}\dot{\neg}U
 =\mathbb{R}\setminus\{0,1\},
\]
ce qui fournit un contre-modèle topologique au tiers exclu.

\subsection{Réalisabilité propositionnelle}

\thingy Reprenons une numérotation
\[
\Phi_e:\mathbb{N}\dashrightarrow\mathbb{N}
\]
des fonctions générales récursives. On interprète maintenant une
formule par un ensemble d'entiers, ses réalisateurs.

Pour $P,Q\subseteq\mathbb{N}$, on pose
\[
P\mathbin{\dottedland}Q
 :=\{\langle m,n\rangle:m\in P,\ n\in Q\},
\]
\[
P\mathbin{\dottedlor}Q
 :=\{\langle0,m\rangle:m\in P\}
   \cup\{\langle1,n\rangle:n\in Q\},
\]
et
\[
P\mathbin{\dottedlimp}Q
 :=\{e\in\mathbb{N}:
     \forall m\in P,\;
     \Phi_e(m)\mathord{\downarrow}\in Q\}.
\]
Enfin,
\[
\dottedtop:=\mathbb{N},
\qquad
\dottedbot:=\varnothing.
\]
Une formule $\varphi(A_1,\ldots,A_r)$ est dite
\defin{réalisable} lorsqu'il existe un même entier $n$ appartenant à
\[
\bigcap_{P_1,\ldots,P_r\subseteq\mathbb{N}}
\dot{\varphi}(P_1,\ldots,P_r).
\]

\begin{thm}[Dragalin, Troelstra, Nelson]
La sémantique de la réalisabilité est correcte pour le calcul
propositionnel intuitionniste.
\end{thm}

Le réalisateur peut être extrait de la preuve : il s'agit d'une autre
forme de la correspondance de Curry--Howard.

Cette sémantique n'est toutefois pas complète. Il existe des formules
réalisables qui ne sont pas démontrables intuitionnistement.

\thingy La formule $A\land B\Rightarrow B\land A$ est réalisée par un
programme qui permute les composantes d'un couple. C'est exactement
l'algorithme correspondant au terme
\[
\lambda z.\langle\pi_2z,\pi_1z\rangle
\]
déjà obtenu par Curry--Howard.

En revanche, $A\lor\neg A$ ne peut pas être uniformément réalisé : il
faudrait produire le même entier qui, pour toute partie
$P\subseteq\mathbb{N}$, indique soit un élément de $P$, soit un
programme certifiant que $P$ est vide.

\subsection{Sémantique des problèmes finis}

\thingy Un \defin{problème fini} est un couple $(X,S)$ où $X$ est un
ensemble fini non vide de candidats et $S\subseteq X$ l'ensemble de ses
solutions.

On définit
\[
(X,S)\mathbin{\dottedland}(Y,T)
 :=(X\times Y,S\times T),
\]
\[
(X,S)\mathbin{\dottedlor}(Y,T)
 :=(X\uplus Y,S\uplus T),
\]
et
\[
(X,S)\mathbin{\dottedlimp}(Y,T)
 :=(Y^X,U),
\]
où
\[
U:=\{f:X\to Y:f(S)\subseteq T\}.
\]
Enfin,
\[
\dottedtop:=(\{\bullet\},\{\bullet\}),
\qquad
\dottedbot:=(\{\bullet\},\varnothing).
\]

Une formule est dite valide au sens de Medvedev lorsqu'il existe, pour
tous les ensembles de candidats choisis, un même candidat qui est une
solution du problème interprétant la formule, quels que soient les
sous-ensembles de solutions.

\begin{thm}[Medvedev]
La sémantique des problèmes finis est correcte pour le calcul
propositionnel intuitionniste.
\end{thm}

Elle n'est pas complète. Elle formalise néanmoins de manière précise
l'idée de Brouwer--Heyting--Kolmogorov selon laquelle une preuve d'une
implication est une fonction transformant les solutions du premier
problème en solutions du second.

\section{Inférence de type à la Hindley--Milner}

\subsection{Désannotation et problème de l'inférence}

\thingy Le $\lambda$CST présenté jusqu'ici porte toutes les annotations
de type. La \defin{désannotation} d'un terme consiste à les supprimer.
Par exemple,
\[
\lambda(x:\alpha).
\lambda(f:(\alpha\to\beta)\to\beta\to\gamma).
f(\lambda(h:\alpha\to\beta).hx)
\]
se désannote en
\[
\lambda xf.\,f(\lambda h.hx).
\]

Le problème de l'inférence de type est alors le suivant :
\begin{itemize}
\item décider si un terme non typé est la désannotation d'un terme du
  $\lambda$CST ;
\item si oui, retrouver ses annotations et son type ;
\item trouver même un type principal, c'est-à-dire le type le plus
  général dont tous les autres s'obtiennent par substitution.
\end{itemize}

L'algorithme de Hindley--Milner, également appelé algorithme de
Damas--Milner ou algorithme W, accomplit ces tâches. Il transforme les
variables de type opaques du $\lambda$CST en un véritable mécanisme de
polymorphisme.

\subsection{Collection et résolution des contraintes}

\thingy La première phase de l'algorithme consiste à collecter des
contraintes. On attribue une variable de type fraîche à chaque
variable et à chaque résultat intermédiaire.

Pour une variable $v$, on renvoie le type qui lui est associé dans le
contexte. Pour une abstraction $\lambda v.e$, on attribue à $v$ un
type frais $\eta$, on infère un type $\zeta$ pour $e$, puis on renvoie
$\eta\to\zeta$. Pour une application $fx$, si les types provisoires de
$f$ et $x$ sont $\zeta$ et $\xi$, on introduit un type frais $\eta$,
on ajoute la contrainte
\[
\zeta=(\xi\to\eta)
\]
et on attribue le type $\eta$ à l'application.

À titre d'exemple, pour
\[
\lambda xyz.xz(yz),
\]
on collecte
\[
x:\eta_1,\qquad y:\eta_2,\qquad z:\eta_3,
\]
\[
xz:\eta_4
\quad\text{avec}\quad
\eta_1=(\eta_3\to\eta_4),
\]
\[
yz:\eta_5
\quad\text{avec}\quad
\eta_2=(\eta_3\to\eta_5),
\]
et
\[
xz(yz):\eta_6
\quad\text{avec}\quad
\eta_4=(\eta_5\to\eta_6).
\]
Le type provisoire renvoyé est
\[
\eta_1\to\eta_2\to\eta_3\to\eta_6.
\]

\thingy La deuxième phase résout les contraintes par
\defin{unification}. Une substitution des variables de type est
cherchée de manière à rendre égaux les deux membres de chaque
contrainte.

L'algorithme procède ainsi :
\begin{itemize}
\item une contrainte déjà satisfaite est supprimée ;
\item si l'un des membres est une variable $\eta$ qui n'apparaît pas
  dans l'autre membre, on substitue cet autre membre à $\eta$ partout ;
\item si les deux membres sont des types fonctions, on les décompose
  composante par composante ;
\item dans tous les autres cas, l'unification échoue.
\end{itemize}
La condition interdisant de remplacer $\eta$ par un type contenant
$\eta$ est appelée \defin{test d'occurrence}. Elle interdit notamment
de résoudre
\[
\eta=(\eta\to\zeta).
\]

\thingy Reprenons
\[
t:=(\lambda xy.x)(\lambda x'y'.x').
\]
La première phase donne les variables de type
\[
x:\eta_1,\quad y:\eta_2,\quad
x':\eta_3,\quad y':\eta_4
\]
et la contrainte
\[
(\eta_1\to\eta_2\to\eta_1)
 =
((\eta_3\to\eta_4\to\eta_3)\to\eta_5).
\]
L'unification produit successivement
\[
\eta_1\mapsto(\eta_3\to\eta_4\to\eta_3)
\]
et
\[
\eta_5\mapsto(\eta_2\to\eta_3\to\eta_4\to\eta_3).
\]
On obtient finalement le typage
\[
\vdash
(\lambda(x:\eta_3\to\eta_4\to\eta_3).
 \lambda(y:\eta_2).x)
(\lambda(x':\eta_3).\lambda(y':\eta_4).x')
:
\eta_2\to\eta_3\to\eta_4\to\eta_3.
\]

\begin{thm}[Propriétés de l'algorithme de Hindley--Milner]
L'algorithme d'unification termine. Lorsqu'il renvoie une substitution,
celle-ci satisfait les contraintes. S'il existe une solution, il
renvoie une solution principale dont toute autre solution s'obtient
par substitution.
\end{thm}

\begin{proof}
La terminaison se montre par une double récurrence, d'abord sur le
nombre de variables de type restant à éliminer, puis sur la taille
totale des types présents dans les contraintes.

La correction se vérifie directement à chaque transformation du
système de contraintes. La propriété de principalité se montre par
induction sur les appels récursifs de l'algorithme : toute solution du
système initial doit effectuer les identifications imposées par
l'étape courante, puis se factorise par la solution renvoyée pour le
système simplifié.

\review{Les diapositives ne donnent que les idées de ces preuves. La
notion exacte de composition des substitutions et la mesure de
terminaison devront être détaillées dans une version ultérieure.}
\end{proof}

\thingy Le terme $\lambda x.xx$ n'est pas typable dans le
$\lambda$CST. La collecte des contraintes imposerait
\[
\eta_1=(\eta_1\to\eta_2),
\]
ce que le test d'occurrence interdit. Cette impossibilité est cohérente
avec la normalisation forte : si $\lambda x.xx$ était typable, le terme
\[
(\lambda x.xx)(\lambda x.xx)
\]
le serait aussi, alors qu'il se réduit indéfiniment à lui-même.

On peut étendre certains systèmes à des types récursifs sans
constructeur, en supprimant ce test. L'option
\texttt{ocaml -rectypes} permet par exemple d'accepter de tels types.
Cette extension doit cependant être utilisée avec précaution.

\subsection{Polymorphisme du \texttt{let} et restriction de valeur}

\thingy Dans un langage fonctionnel,
\[
\texttt{let }v=x\texttt{ in }t
\]
peut être vu comme un sucre syntaxique pour
\[
(\lambda v.t)x.
\]
Cette traduction pose toutefois un problème pour le polymorphisme. Si
une même valeur doit être utilisée à plusieurs types dans $t$, le
typage ordinaire d'une abstraction impose un type unique à $v$.

L'idée de Hindley--Milner est de typer d'abord $x$, puis de
\defin{généraliser} les variables de type qui n'apparaissent pas dans
le contexte extérieur. Chaque occurrence de $v$ dans $t$ reçoit
ensuite des variables de type fraîches obtenues par instanciation du
type généralisé.

Cette opération permet par exemple d'utiliser une même fonction
\texttt{twice} aussi bien sur des entiers que sur des booléens dans le
corps d'un même \texttt{let}.

\thingy Lorsque le langage possède des effets de bord et des références
mutables, une généralisation sans restriction serait dangereuse.
Considérons schématiquement
\[
\texttt{let r = ref (fun x -> x) in
  r := (fun x -> x+1); (!r) true}.
\]
Si le type de \texttt{r} était généralisé, la même cellule pourrait être
vue à la fois comme une référence vers une fonction
$\texttt{int}\to\texttt{int}$ et comme une référence vers une fonction
$\texttt{bool}\to\texttt{bool}$. Le programme appliquerait alors une
fonction entière à un booléen.

OCaml utilise donc la \defin{restriction de valeur} : le membre droit
d'un \texttt{let} n'est pleinement généralisé que lorsqu'il est une
valeur statique déterminable à la compilation. Dans un langage
purement fonctionnel comme Haskell, ce problème ne se pose pas de la
même manière, puisque les valeurs et les fonctions sont pures.